% Bagian: first-order-logic
% Bab: models-theories
% Subbagian: set-theory

\documentclass[../../../include/open-logic-section]{subfiles}

\begin{document}

\olfileid[id]{fol}{mat}{set}

\olsection{Teori Himpunan}

Hampir seluruh matematika dapat dikembangkan di dalam teori himpunan.
Mengembangkan matematika dalam teori ini melibatkan sejumlah hal. Pertama,
diperlukan suatu himpunan aksioma bagi relasi~$\in$. Sejumlah sistem aksioma
yang berbeda telah dikembangkan, kadang dengan sifat-sifat~$\in$ yang saling
bertentangan. Sistem aksioma yang dikenal sebagai $\Log{ZFC}$, yaitu teori
himpunan Zermelo--Fraenkel dengan aksioma pilihan, merupakan yang paling
menonjol: sistem itu sejauh ini paling luas digunakan dan dikaji karena
ternyata aksioma-aksiomanya cukup untuk membuktikan hampir semua hal yang
diharapkan dapat dibuktikan oleh matematikawan. Namun, sebelum hal itu dapat
dipastikan, pertama-tama perlu dijernihkan bagaimana kita bahkan dapat
\emph{mengungkapkan} semua hal yang ingin diungkapkan oleh matematikawan.
Sebagai permulaan, bahasa tersebut tidak memuat !!{constant}s atau
!!{function}s, sehingga sepintas tidak jelas bahwa kita dapat berbicara
tentang himpunan tertentu (seperti $\emptyset$ atau $\Nat$), tentang operasi
pada himpunan (seperti $X \cup Y$ dan $\Pow{X}$), apalagi tentang konstruksi
lain yang melibatkan hal-hal selain himpunan, seperti relasi dan fungsi.

Sebagai langkah awal, ``merupakan anggota dari'' bukan satu-satunya relasi
yang kita minati: ``merupakan himpunan bagian dari'' tampak hampir sama
pentingnya. Namun, kita dapat \emph{mendefinisikan} ``merupakan himpunan
bagian dari'' berdasarkan ``merupakan anggota dari.'' Untuk melakukannya, kita
harus menemukan !!a{formula}~$!A(x, y)$ dalam bahasa teori himpunan yang
dipenuhi oleh suatu pasangan himpunan~$\tuple{X, Y}$ jika dan hanya jika
$X \subseteq Y$. Namun, $X$ merupakan himpunan bagian dari $Y$ tepat jika
semua !!{element}s dari~$X$ juga merupakan !!{element}s dari~$Y$. Jadi, kita
dapat mendefinisikan $\subseteq$ dengan formula
\[
\lforall[z][(z \in x \lif z \in y)]
\]
Sekarang, setiap kali hendak menggunakan relasi~$\subseteq$ dalam suatu
formula, kita dapat menggantinya dengan formula tersebut (dengan $x$ dan $y$
diganti sebagaimana mestinya, dan variabel terikat~$z$ dinamai ulang jika
perlu). Misalnya, ekstensionalitas himpunan berarti bahwa jika sebarang
himpunan~$x$ dan $y$ termuat satu sama lain, maka $x$ dan $y$ harus merupakan
himpunan yang sama. Hal ini dapat diungkapkan dengan
$\lforall[x][\lforall[y][((x \subseteq y \land y \subseteq x) \lif x =
y)]]$, atau, jika kita mengganti $\subseteq$ dengan definisi di atas, dengan
\[
\lforall[x][\lforall[y][((\lforall[z][(z \in x \lif z \in y)] \land
    \lforall[z][(z \in y \lif z \in x)]) \lif x = y)]].
\]
Ini sebenarnya merupakan salah satu aksioma $\Log{ZFC}$, yakni ``aksioma
ekstensionalitas.''

Tidak ada !!{constant} bagi $\emptyset$, tetapi kita dapat mengungkapkan
``$x$ kosong'' dengan $\lnot \lexists[y][y \in x]$. Kemudian
``$\emptyset$ ada'' menjadi !!{sentence}~$\lexists[x][\lnot
\lexists[y][y \in x]]$. Ini merupakan aksioma lain dari~$\Log{ZFC}$.
(Perhatikan bahwa aksioma ekstensionalitas menyiratkan bahwa hanya ada satu
himpunan kosong.) Setiap kali hendak berbicara tentang $\emptyset$ dalam
bahasa teori himpunan, kita akan menuliskannya sebagai ``terdapat suatu
himpunan yang kosong dan \dots'' Sebagai contoh, untuk mengungkapkan fakta
bahwa $\emptyset$ merupakan himpunan bagian dari setiap himpunan, kita dapat
menulis
\[
\lexists[x][(\lnot\lexists[y][y \in x] \land \lforall[z][x \subseteq
    z])]
\]
yang, tentu saja, pada gilirannya mengharuskan $x \subseteq z$ diganti dengan
definisinya.

Untuk berbicara tentang operasi pada himpunan, seperti $X \cup Y$ dan
$\Pow{X}$, kita harus menggunakan siasat serupa. Tidak ada simbol fungsi dalam
bahasa teori himpunan, tetapi kita dapat mengungkapkan relasi fungsional
$X \cup Y = Z$ dan $\Pow{X} = Y$ dengan
\begin{align*}
& \lforall[u][((u \in x \lor u \in y) \liff u \in z)]\\
& \lforall[u][(u \subseteq x \liff u \in y )]
\end{align*}
karena !!{element}s dari $X \cup Y$ tepat merupakan himpunan-himpunan yang
menjadi !!{element}s dari~$X$ atau !!{element}s dari~$Y$, dan !!{element}s
dari $\Pow{X}$ tepat merupakan himpunan-himpunan bagian dari~$X$. Namun, hal
ini tidak memungkinkan kita menggunakan $x \cup y$ atau $\Pow{x}$ seolah-olah
keduanya merupakan suku: kita hanya dapat menggunakan keseluruhan !!{formula}s
yang mendefinisikan relasi $X \cup Y = Z$ dan $\Pow{X} = Y$. Bahkan, kita
belum mengetahui bahwa relasi-relasi ini pernah terpenuhi, yakni kita belum
mengetahui bahwa gabungan dan himpunan kuasa selalu ada. Misalnya,
!!{sentence} $\lforall[x][\lexists[y][\Pow{x} = y]]$ merupakan aksioma lain
dari~$\Log{ZFC}$ (aksioma himpunan kuasa).

Lalu bagaimana dengan pembicaraan mengenai pasangan terurut atau fungsi? Di
sini kita harus menjelaskan bagaimana kita dapat memandang pasangan terurut
dan fungsi sebagai jenis himpunan khusus. Salah satu cara mendefinisikan
pasangan terurut $\tuple{x, y}$ ialah sebagai himpunan
$\{\{x\}, \{x, y\}\}$. Namun, seperti sebelumnya, kita tidak dapat
memperkenalkan !!a{function} yang menamai himpunan ini; kita hanya dapat
mendefinisikan relasi $\tuple{x, y} = z$, yakni
$\{\{x\}, \{x, y\}\} = z$:
\[
\lforall[u][(u \in z \liff (\lforall[v][(v \in u \liff v = x)] \lor
  \lforall[v][(v \in u \liff (v = x \lor v = y))]))]
\]
Formula ini menyatakan bahwa !!{element}s~$u$ dari~$z$ tepat merupakan
himpunan-himpunan yang mempunyai $x$ sebagai satu-satunya !!{element}, atau
mempunyai $x$ dan~$y$ sebagai satu-satunya !!{element}s (dengan kata lain,
himpunan yang identik dengan $\{x\}$ atau identik dengan $\{x, y\}$). Setelah
mempunyai hal ini, kita dapat menyatakan hal-hal lebih lanjut, misalnya bahwa
$X \times Y = Z$:
\[
\lforall[z][(z \in Z \liff \lexists[x][\lexists[y][(x \in
        X \land y \in Y \land \tuple{x, y} = z)]])]
\]

Suatu fungsi $f \colon X \to Y$ dapat dipandang sebagai relasi $f(x) = y$,
yakni sebagai himpunan pasangan~$\Setabs{\tuple{x,y}}{f(x) = y}$. Kemudian
kita dapat mengatakan bahwa suatu himpunan~$f$ merupakan fungsi dari $X$ ke
$Y$ jika (a) himpunan itu merupakan relasi $\subseteq X \times Y$, (b)
himpunan itu total, yakni bagi semua $x \in X$ terdapat suatu $y \in Y$
sedemikian sehingga $\tuple{x, y} \in f$, dan (c) himpunan itu fungsional,
yakni setiap kali $\tuple{x, y}, \tuple{x, y'} \in f$, berlaku $y = y'$
(karena nilai fungsi harus tunggal). Jadi, ``$f$ merupakan fungsi dari $X$ ke
$Y$'' dapat ditulis sebagai:
\begin{align*}
 \lforall[u][(u \in f \lif {}] & \lexists[x][\lexists[y][(x \in X \land y \in
      Y \land \tuple{x, y} = u)]]) \land {}\\
 \lforall[x][(x \in X \lif {}] &
      (\lexists[y][(y \in Y \land \mathrm{maps}(f, x, y))] \land {}\\
& (\lforall[y][\lforall[y'][((\mathrm{maps}(f, x, y) \land
    \mathrm{maps}(f, x, y')) \lif y = y')]])))
\end{align*}
dengan $\mathrm{maps}(f, x, y)$ menyingkat $\lexists[v][(v \in f \land
  \tuple{x, y} = v)]$ (!!{formula} ini mengungkapkan ``$f(x) = y$'').

Sekarang juga tidak sukar untuk mengungkapkan bahwa $f\colon X \to Y$
bersifat !!{injective}, misalnya:
\begin{multline*}
 f \colon X \to Y \land \lforall[x][\lforall[x'][((x \in X \land x' \in
     X \land {}]] \\
 \lexists[y][(\mathrm{maps}(f, x, y) \land \mathrm{maps}(f,
        x', y))]) \lif x = x')
\end{multline*}
Suatu fungsi~$f\colon X \to Y$ bersifat !!{injective} jika dan hanya jika,
setiap kali $f$ memetakan $x, x' \in X$ ke satu~$y$, berlaku $x = x'$.
Jika kita menyingkat formula ini sebagai $\mathrm{inj}(f, X, Y)$, kita sudah
dapat menyatakan sesuatu yang tidak sepele seperti teorema Cantor dalam bahasa
teori himpunan: tidak ada fungsi !!{injective} dari $\Pow{X}$ ke $X$:
\[
\lforall[X][\lforall[Y][(\Pow{X} = Y \lif
    \lnot\lexists[f][\mathrm{inj}(f, Y, X)])]]
\]

Orang mungkin menduga bahwa teori himpunan memerlukan aksioma lain yang
menjamin keberadaan suatu himpunan bagi setiap sifat pendefinisi. Jika $!A(x)$
merupakan formula teori himpunan dengan variabel~$x$ bebas, kita dapat
meninjau !!{sentence}
\[
\lexists[y][\lforall[x][(x \in y \liff !A(x))]].
\]
!!^{sentence} ini menyatakan bahwa terdapat suatu himpunan~$y$ yang
!!{element}snya tepat semua $x$ yang memenuhi~$!A(x)$. Skema ini disebut
``prinsip komprehensi.'' Prinsip ini tampak sangat berguna; sayangnya, prinsip
ini tidak konsisten. Ambil $!A(x) \ident \lnot x \in x$, maka prinsip
komprehensi menyatakan
\[
\lexists[y][\lforall[x][(x \in y \liff x \notin x)]],
\]
yakni prinsip itu menyatakan adanya himpunan semua himpunan yang bukan
!!{element}s dari dirinya sendiri. Himpunan seperti itu tidak mungkin
ada---inilah Paradoks Russell. $\Log{ZFC}$ sebenarnya memuat versi terbatas
---dan konsisten---dari prinsip ini, yakni prinsip separasi:
\[
\lforall[z][\lexists[y][\lforall[x][(x \in y \liff (x \in z \land
      !A(x)))]]].
\]

\begin{prob}
Tunjukkan bahwa prinsip komprehensi tidak konsisten dengan memberikan
!!a{derivation} yang menunjukkan
\[
\lexists[y][\lforall[x][(x \in y \liff x \notin x)]] \Proves \lfalse.
\]
Mungkin akan membantu jika Anda terlebih dahulu menunjukkan
$(A \lif \lnot A) \land (\lnot A \lif A) \Proves \lfalse$.
\end{prob}
\end{document}
